\documentclass[10pt]{article}
\usepackage{amsmath,amssymb,geometry,url}
\geometry{margin=1in}
\setlength{\emergencystretch}{3em}
\title{Prime-power groups with identical subgroup lattices\\A disproof of conjecture 00000004294}
\author{gaochengzhecpu (AI-assisted submission)}
\date{3 October 2026}
\begin{document}
\maketitle

\paragraph{The assertion being disproved.}
The conjecture asserts, in particular, that the subgroup lattice is a complete invariant on the class of abelian groups of prime-power order. Thus, for two groups in this class, isomorphic subgroup lattices would imply isomorphic groups. We disprove precisely this necessary part of the statement.

The prime is allowed to vary: the English statement says ``the class of prime-power-order groups,'' and the Chinese statement uses the same prime-power wording. This is not the different assertion that one first fixes a common prime $p$ and then compares only finite abelian $p$-groups. The example below does not disprove that restricted assertion. No claim about the other, rank-one torsion-free, part of the proposed classification is needed.

\paragraph{Counterexample.}
Let
\[
 G=\mathbb Z/2\mathbb Z,\qquad H=\mathbb Z/3\mathbb Z
\]
with their additive group structures. Both are abelian and have prime-power order, namely $2^1$ and $3^1$.

The only subgroups of $G$ are $\{0\}$ and $G$: a subgroup containing its only nonzero element $1$ is the whole group. The only subgroups of $H$ are likewise $\{0\}$ and $H$. Indeed, a subgroup containing $1$ contains $1+1=2$, and one containing $2$ contains $2+2=1$ modulo $3$.

Consequently their entire subgroup lattices are the same two-element chain, and the map
\[
 \{0\}\longmapsto\{0\},\qquad G\longmapsto H
\]
is a lattice isomorphism. It preserves both inclusion and the lattice operations: the meet is subgroup intersection, and the join is the subgroup generated by the union. On the other hand, $G$ and $H$ cannot be isomorphic because their cardinalities are respectively $2$ and $3$. Thus the subgroup lattice is not a complete invariant on the stated class. This contradicts the conjecture.

\newpage
\paragraph{Exact scope of the formalization.}
The accompanying \texttt{lean/Main.lean} is a self-contained Lean 4.19.0 project importing only \texttt{Std}. The types \texttt{Fin 2} and \texttt{Fin 3} carry the actual addition modulo $2$ and $3$. All abelian group laws are checked. Primality is the usual positive-divisor condition, quantified over all natural numbers; the orders are certified by bijections with \texttt{Fin} of the relevant prime power.

A \texttt{Subgroup} is an arbitrary proposition-valued membership predicate satisfying closure under zero, addition, and additive inverses. The formalization therefore quantifies over all actual subgroups; it does not rely on a selection of bit masks or on an unproved enumeration. The lemmas \path{subgroup2_membership} and \path{subgroup3_membership} prove, for every such subgroup and every group element, the exact formula
\[
 x\in K\quad\Longleftrightarrow\quad x=0\ \text{or}\ 1\in K.
\]
Using this classification, \path{subgroupLatticesIso} gives explicit mutually inverse maps on the entire subgroup types and proves that they preserve and reflect inclusion. The lattice meet is defined by predicate intersection; the join is defined as the intersection of all subgroups containing both inputs. Generic proofs show that the order isomorphism preserves these operations, and the specialized theorems \path{lattice_meet_preserved} and \path{lattice_join_preserved} apply them to the present example.

A \texttt{GroupIso} includes an additive map and a two-sided inverse. The theorem \path{groups_not_isomorphic} rules out every such isomorphism by showing that no injection from \texttt{Fin 3} to \texttt{Fin 2} exists. The final theorem \path{conjecture04294_false} negates the universal completeness implication on the prime-power-order class. It uses both the actual subgroup-lattice isomorphism and the nonisomorphism of the groups.

All finite checks are kernel-checked by \texttt{decide}. The printed axiom lists for the final theorem contain only Lean's standard \texttt{propext} and \texttt{Quot.sound}; there are no proof placeholders, external decision oracles, or additional axioms. The project uses \texttt{warningAsError=true} and is built by \texttt{lake build}.
\end{document}
